\documentclass[11pt,a4paper]{article}
\usepackage[utf8]{inputenc}
\usepackage{amsmath,amssymb}
\usepackage{graphicx}
\usepackage{booktabs}
\usepackage{hyperref}
\usepackage{cite}
\usepackage{listings}
\usepackage{xcolor}
\usepackage{geometry}
\geometry{margin=1in}

\lstset{
    basicstyle=\ttfamily\footnotesize,
    keywordstyle=\color{blue},
    stringstyle=\color{red},
    commentstyle=\color{gray},
    breaklines=true,
    frame=single
}

\title{\textbf{AI Control Plane:\\A First-Principles Architectural Specification for Enterprise AI Governance, Observability, and Operational Control}}

\author{
    \textbf{Larry Arnold}\\
    Founder, LA Builds \& Systems Architect\\
    \texttt{arnoldlarry15@github.com}
}
\date{January 2026 (Revised: October 2026)}

\begin{document}

\maketitle

\begin{abstract}
As enterprises transition from single-prompt experiments to autonomous multi-agent systems, organizations face an urgent architectural challenge: AI systems operate with opaque execution paths, unvetted tool access, and zero runtime policy enforcement. We specify the \textbf{AI Control Plane}, an infrastructure-layer control architecture that strictly decouples the AI \textit{compute plane} (model inference, fine-tuning, hosted LLM APIs) from the AI \textit{control plane} (identity, declarative guardrails, telemetry, and kill-switch mediation). Built upon a choke-point execution gateway, machine-enforced declarative policies, immutable audit logging, and hardware-level kill-switch isolation, this specification provides deterministic oversight without degrading developer velocity. We define formal system components, threat boundary models, state transition lifecycles, and a reference Python SDK implementation compatible with NIST AI RMF, OWASP LLM Top 10, and EU AI Act compliance standards.
\end{abstract}

\vspace{1em}
\noindent \textbf{Keywords:} AI Governance, Control Plane Architecture, Choke-Point Gateway, Declarative Guardrails, Enterprise Observability, Kill-Switch Mechanics.

\section{Introduction and Strategic Context}
Enterprise software architectures have repeatedly demonstrated that when computational capabilities expand rapidly, governance and security must be decoupled from runtime computation. The cloud revolution separated infrastructure control planes (Kubernetes, AWS IAM) from user compute workloads. 

Currently, enterprise adoption of Large Language Models (LLMs) and autonomous agents mimics early unmanaged cloud deployments: individual engineering teams embed direct API calls across disparate services, creating fragmented ``shadow AI'' footprints without centralized inventory, permission auditing, or immediate fail-safe capabilities.

The \textbf{AI Control Plane} resolves this structural deficit. Rather than competing with foundational model providers or attempting to train proprietary models, the Control Plane inserts a deterministic, low-latency governance layer between applications and AI inference endpoints.

\section{Architectural First Principles}
The AI Control Plane is governed by five non-negotiable architectural axioms:
\begin{enumerate}
    \item \textbf{Decoupled Control vs. Compute:} The Control Plane does not perform inference; it regulates, audits, and mediates access to compute.
    \item \textbf{Identity Precedes Execution:} If an AI model, agent, or workflow is not explicitly cataloged in the Registry, all downstream API calls are rejected.
    \item \textbf{Declarative Machine-Enforced Policy:} Rules beat heuristics. Policies are version-controlled, human-auditable YAML/JSON artifacts evaluated dynamically prior to inference.
    \item \textbf{Unified Gateway Choke-Point:} All ingress prompts, tool invocations, and egress responses pass through an intercepting gateway. Bypassing the gateway results in network-level disconnection.
    \item \textbf{Deterministic Fail-Safe and Kill-Switch:} System administrators possess instantaneous, cryptographically verified kill-switch authority at the individual agent, team, and global organizational tier.
\end{enumerate}

\section{Core System Pillars}

\subsection{1. AI Asset Registry}
The Registry serves as the centralized source of truth for all autonomous and semi-autonomous AI artifacts in the organization. Every entity is registered with:
\begin{itemize}
    \item \textbf{AI Entity ID:} Unique cryptographic identifier (UUIDv4).
    \item \textbf{Entity Classification:} Static Model Wrapper, Autonomous Tool-Calling Agent, Multi-Agent Swarm, or Scheduled Daemon.
    \item \textbf{Model Provider \& Version:} Explicit upstream endpoint (e.g., \texttt{claude-3-5-sonnet}, \texttt{gpt-4o}).
    \item \textbf{Risk Tier:} Tier 1 (Public Read-Only), Tier 2 (Internal Data Read), Tier 3 (Write / Transactional Capability), Tier 4 (Autonomous External Action).
    \item \textbf{Owner \& Cost Center:} Designated engineering owner, team accountability, and hard budget limits.
\end{itemize}

\subsection{2. Declarative Policy \& Guardrail Engine}
Policies are expressed declaratively, allowing security and legal teams to review and version guardrails via standard Git workflows:

\begin{lstlisting}[language=SQL]
policy_version: "2026.1"
agent_id: "agent_support_triage_042"
risk_tier: 3
allowed_tools:
  - "read_knowledge_base"
  - "draft_ticket_reply"
forbidden_tools:
  - "execute_refund"
  - "modify_database_record"
data_boundaries:
  pii_redaction: true
  allowed_tenants: ["tenant_us_east"]
on_violation:
  action: "BLOCK_AND_ESCALATE"
  notify: "secops-alert@enterprise.internal"
\end{lstlisting}

\subsection{3. The Choke-Point Execution Gateway}
The gateway intercepts all inbound and outbound AI traffic. For each request:
\begin{equation}
\mathcal{G}(R) = \mathcal{E}_{\text{policy}}(\mathcal{I}(R)) \to \mathcal{M}_{\text{inference}}(R) \to \mathcal{E}_{\text{output}}(O) \to \mathcal{L}_{\text{audit}}(E)
\end{equation}
If policy evaluation $\mathcal{E}_{\text{policy}}$ detects an authorization breach, the gateway aborts execution immediately with a standardized error code (\texttt{403\_AI\_POLICY\_VIOLATION}) and logs the incident.

\subsection{4. Immutable Observability and Audit Layer}
Every interaction produces an immutable telemetry record storing:
\begin{itemize}
    \item Ingress prompt hash and sanitized content
    \item Evaluated policy rules and millisecond latency overhead
    \item Raw upstream model latency, token consumption, and cost
    \item Tool-call invocations and argument payloads
    \item Full temporal replay capability for post-incident forensics
\end{itemize}

\subsection{5. Hierarchical Kill-Switch Architecture}
The Control Plane provides sub-second revocation across three operational scopes:
\begin{itemize}
    \item \textbf{Targeted Agent Kill-Switch:} Suspends a malfunctioning agent without impacting unrelated organizational services.
    \item \textbf{Tool Revocation Switch:} Instantly disarms specific tool capabilities (e.g., file-system writes or database updates) while leaving read-only operations active.
    \item \textbf{Global Air-Gap Switch:} Severing all external LLM provider API connections during severe active compromise.
\end{itemize}

\section{Threat Boundary Model}

\begin{table}[htbp]
\centering
\caption{Threat Vector Mitigation Across the AI Control Plane}
\label{tab:threat_model}
\begin{tabular}{lll}
\toprule
\textbf{Threat Vector} & \textbf{Risk Level} & \textbf{Control Plane Mitigation Mechanism} \\
\midrule
Indirect Prompt Injection & Critical & Gateway content pre-filtering \& tool argument validation \\
Data Exfiltration via Tool Calls & High & Declarative egress policy \& destination whitelisting \\
Autonomous Financial Runaway & High & Hard token budget quotas \& human-in-the-loop limits \\
Unvetted Shadow Agents & Medium & Zero-trust registry enforcement (unregistered calls rejected) \\
Post-Breach Tampering & High & Write-once append-only cryptographic audit logging \\
\bottomrule
\end{tabular}
\end{table}

\section{Reference Software Architecture}
The standard reference implementation organizes into modular microservices:
\begin{itemize}
    \item \texttt{control_plane/registry}: Entity lifecycle and risk-tier registry.
    \item \texttt{control_plane/gateway}: Low-latency asynchronous HTTP/gRPC reverse proxy.
    \item \texttt{control_plane/policy}: Abstract syntax tree policy compiler and evaluator.
    \item \texttt{control_plane/kill_switch}: Distributed Redis/eBPF circuit breaker.
    \item \texttt{control_plane/observability}: OpenTelemetry-compatible tracing and audit sink.
    \item \texttt{sdk/python}: Drop-in client library requiring $< 5$ lines of configuration.
\end{itemize}

\section{Regulatory and Compliance Alignment}
The architecture directly satisfies regulatory mandates:
\begin{itemize}
    \item \textbf{NIST AI RMF 1.0:} Fulfills Govern (1.1, 1.2), Map (1.5), Measure (2.6), and Manage (2.3) requirements through centralized inventory and telemetry.
    \item \textbf{EU AI Act:} Implements Article 14 (Human Oversight) and Article 12 (Record-Keeping) for high-risk autonomous systems.
    \item \textbf{SOC2 Type II:} Provides tamper-proof access control evidence and automated change logging.
\end{itemize}

\section{Conclusion}
The AI Control Plane delivers the necessary architectural maturity for modern AI engineering. By separating operational control from underlying inference compute, organizations obtain deterministic safety, continuous auditability, and immediate incident mitigation without constraining rapid model adoption.

\bibliographystyle{ieeetr}
\bibliography{references}

\end{document}
